\section{Predication as the seed out of potential}

\begin{proposition}[Entity and predicate]\label{thm:entity-predicate}\lean{Obsmetrologia.Entity}\leanok\uses{def:mark, thm:kind}
An \emph{entity} is a mark in identity-carrying role, and a
\emph{predicate} is a mark in content-carrying role. At this base both are
abbreviations for \texttt{Mark}; typed predicate structure enters later
through positing and valuation.
\end{proposition}

\begin{proposition}[The predication cell]\label{thm:predication}\lean{Obsmetrologia.PredicationCell}\leanok\uses{thm:entity-predicate, thm:articulation-cell-corner}
The predication cell is the oriented triple $(s,p,o)$: subject entity,
predicate mark, and object entity. Orientation belongs to the datum, so the
same predicate read from opposite endpoints is a different cell off the
diagonal.
\end{proposition}

\begin{proposition}[The claim point]\label{thm:claim-point}\lean{Obsmetrologia.Claim}\leanok\uses{thm:predication}
A claim point is the predication cell read as the generated point at which later
value, settlement, and descent can be carried. It is not yet a site object, a
cover, or a glued documentation carrier; those structures are generated later.
\end{proposition}

\begin{proposition}[The reflexive seed]\label{thm:predication-seed}\lean{Obsmetrologia.PredicationCell.seed}\leanok\uses{thm:predication, def:utterance}
The seed of an utterance with token $t$ and claim $c$ is the reflexive
predication cell $(t,c,t)$. Predicating first creates the diagonal assertion
position before later constructions move off the diagonal.
\end{proposition}

\begin{theorem}[The claim seed records the claim-side mark]\label{thm:claim-seed-predicate}\lean{Obsmetrologia.Claim.seed_predicate}\leanok\uses{thm:claim-point, thm:predication-seed}
The claim point generated by an utterance keeps the utterance's claim-side mark
as its predicate coordinate. This is the formal reason the name \emph{Claim}
belongs at the predication-point layer rather than at the later site-cover layer.
\end{theorem}

\begin{proposition}[Predicating proper]\label{thm:predicating-proper}\lean{Obsmetrologia.PredicationCell.predicating}\leanok\uses{thm:predication-seed, thm:articulation-cell-initial}
Predicating is defined on an articulation cell only with a proof that the
cell is in the potential corner. Under that incidence condition it returns
the seed of the cell's utterance.
\end{proposition}
